\documentclass[../../main.tex]{subfiles}
\begin{document}

\begin{center}
  {\fontsize{48pt}{58pt}\selectfont \textbf{2006}\normalsize\textbf{年度}}
\end{center}

\begin{table}[htb]
  \centering
  \caption{出題テーマと難易度}
  \begin{tabular}{|c|p{9cm}|c|}
    \hline
    \textbf{問題番号} & \textbf{テーマ} & \textbf{難易度} \\
    \hline
    第1問 & 整式の平方の整除性と重解 & \\
    \hline
    第2問 & 空間の2線分の交点 & \\
    \hline
    第3問 & 点対称な放物線と接線で囲む面積 & \\
    \hline
    第4問 & 同時に素数となる整数の決定 & \\
    \hline
    第5問 & 三角形の辺上を動く点の重心の範囲 & \\
    \hline
    第6問 & 三角関数を含む定積分の最大値 & \\
    \hline
  \end{tabular}
\end{table}

\begin{table}[htb]
  \centering
  \caption{解答}
  \renewcommand{\arraystretch}{1.4}
  \begin{tabular}{|c|c|p{9cm}|}
    \hline
    \textbf{問題番号} & \textbf{小問} & \textbf{答え} \\
    \hline
    第1問 & -- & $Q(x)=0$は重解をもつ（証明問題） \\
    \hline
    第2問 & -- & $t=-1$のとき交点$\left(\dfrac{4}{3},\dfrac{7}{3},\dfrac{2}{3}\right)$ \\
    \hline
    第3問 & -- & $\dfrac{6+4\sqrt{2}}{3}$ \\
    \hline
    第4問 & -- & $n=3$のみ \\
    \hline
    第5問 & -- & 各辺の三等分点6点を頂点とする六角形の内部（境界を含まない） \\
    \hline
    第6問 & -- & $\theta=\dfrac{\pi}{2}-\alpha$で最大値$\dfrac{\pi}{2}-\alpha-\cos\alpha$ \\
    \hline
  \end{tabular}
\end{table}

\paragraph{試験形態}
6問．試験時間は要確認．

\end{document}
